\chapter*{How to use this guide}
\addcontentsline{toc}{chapter}{How to use this guide}
\markboth{How to use this guide}{How to use this guide}

This document is a self-contained study companion for the MSc course \textbf{AE4262 -- Combustion for Propulsion and Power Technologies} taught at TU Delft. Its aim is that a reader can learn the material from this text alone, without the slides open: the bullet points of the lectures have been turned into continuous prose, and the derivations, definitions and physical reasoning that the slides assume or only sketch have been filled in from the prescribed textbooks.

The backbone and scope follow the AE4262 lecture set~\cite{lecture1}; the depth of the derivations and definitions is drawn from the course reference books, principally Warnatz, Maas \& Dibble~\cite{warnatz2006combustion}, the companion lecture notes by Maas~\cite{maas_combustion_lecture}, Glassman \& Yetter~\cite{glassman2014combustion}, and the lecture notes for the introductory lectures by Langella~\cite{Langella_LectureNotes}, supplemented by Kuo~\cite{kuo2005conservation}, Turns~\cite{turns2000introduction}, Peters~\cite{peters2000turbulent}, Pope~\cite{pope2000turbulent}, the review of Veynante \& Vervisch~\cite{veynante2002turbulent} and the stretched-flame review of Law~\cite{Law1988}.

The material is organised by concept rather than by lecture order. After an introduction that motivates why combustion is difficult and classifies the phenomena (Chapter~\ref{ch:intro}), the chemistry-free thermodynamic description is developed first (Chapter~\ref{ch:thermo}), then finite-rate chemistry (Chapter~\ref{ch:kinetics}) and the governing transport equations (Chapter~\ref{ch:governing}). These tools are applied to laminar premixed (Chapter~\ref{ch:premixed}) and non-premixed (Chapter~\ref{ch:nonpremixed}) flames, to flame stretch and the intrinsic flame instabilities (Chapter~\ref{ch:stretch}), and then extended to turbulence (Chapter~\ref{ch:turbulence}) and turbulent combustion modelling (Chapter~\ref{ch:turbcomb}). The last four chapters treat the experimental side (Chapter~\ref{ch:diagnostics}) and the practical engineering of propulsion combustion: pollutant formation (Chapter~\ref{ch:emissions}), combustor design (Chapter~\ref{ch:combustors}) and hydrogen combustion (Chapter~\ref{ch:hydrogen}).

Citations are given throughout. References of the form ``Lecture~$n$'' point to the AE4262 slide set for that lecture; book and paper citations point to the specific chapters listed in the course guidelines. Where a figure is reproduced from a lecture slide or a textbook, the source is named in the caption.
